\chapter{Về Thái Độ Sống}
\markboth{Về Thái Độ Sống}{On Attitude in Life}

\section{Thông minh mà kiêu ngạo thì không đi được xa.}
\begin{truyen}{Kiêu ngạo không đi xa}{Arrogance Doesn't Go Far}
\chuhoa{C}{ó em học rất giỏi nhưng hay chê bạn chậm, tỏ ra mình hiểu hết. Thầy gọi em ra: ``Thông minh là tài. Nhưng kiêu ngạo là đóng cửa. Khi em cho rằng mình hơn hết, em sẽ không còn muốn học thêm — và người không học thêm sẽ bị người khiêm tốn vượt qua. Thông minh mà kiêu thì không đi được xa.''}
\textit{(A very bright student often looked down on slower classmates. The teacher said: ``Smarts are a gift. But arrogance is a door you shut. When you think you're above everyone, you stop wanting to learn — and those who stop learning get passed by the humble. Smart but arrogant doesn't go far.'')}
\end{truyen}

\begin{baihoc}
Khiêm tốn mở cửa học thêm; kiêu ngạo đóng cửa và đứng yên.
\textit{(Humility opens the door to keep learning; arrogance closes it and you stay still.)}
\end{baihoc}

\begin{ghinhoanh}
\textbf{Short English to remember:}

Smart but arrogant won't get you far.

\end{ghinhoanh}

\begin{tuhocnhanh}
1. Vì sao kiêu ngạo lại hại chính người giỏi? \textit{Why does arrogance hurt even the strong?}

2. Em có khi nào tỏ ra mình hơn bạn không? \textit{Do you ever act as if you're above others?}
\end{tuhocnhanh}
\ngancach

\section{Người tử tế luôn được nhớ lâu hơn người giỏi.}
\begin{truyen}{Tử tế được nhớ lâu}{Kindness Is Remembered Longer}
\chuhoa{T}{hầy hỏi cả lớp: ``Các em nhớ nhất bạn nào trong lớp cũ? Thường không phải bạn giỏi nhất, mà bạn hay giúp em, hay chia sẻ, hay không chê ai. Người giỏi có thể làm em nể; người tử tế làm em nhớ và muốn gần. Hãy cố vừa học vừa tử tế.''}
\textit{(The teacher asked: ``Who do you remember most from your old class? Usually not the top student, but the one who helped, shared, didn't put anyone down. The smart one may impress you; the kind one you remember and want to be near. Try to be both learned and kind.'')}
\end{truyen}

\begin{baihoc}
Giỏi khiến người ta nể; tử tế khiến người ta nhớ và tin.
\textit{(Being good makes people respect you; being kind makes them remember and trust you.)}
\end{baihoc}

\begin{ghinhoanh}
\textbf{Short English to remember:}

Kind people are remembered longer than clever ones.

\end{ghinhoanh}

\begin{tuhocnhanh}
1. Em nhớ nhất ai trong lớp — người giỏi nhất hay người tử tế nhất? \textit{Who do you remember most — the smartest or the kindest?}

2. Em có đang vừa học vừa tử tế không? \textit{Are you both learning and being kind?}
\end{tuhocnhanh}
\ngancach

\section{Đừng chê bạn yếu hơn mình, vì mỗi người giỏi một cách khác nhau.}
\begin{truyen}{Mỗi người giỏi một cách}{Everyone Is Good in Different Ways}
\chuhoa{M}{ột em hay nói ``bạn A dốt toán'', ``bạn B viết xấu''. Thầy nói: ``Em giỏi toán, bạn A có thể giỏi vẽ hoặc giỏi nói chuyện. Mỗi người giỏi một cách. Chê bạn yếu hơn mình ở chỗ này không làm em cao hơn — chỉ làm em hẹp đi. Người lớn thật sự biết nhìn ra điểm mạnh của người khác.''}
\textit{(A student often said ``A is bad at math'', ``B writes ugly''. The teacher said: ``You're good at math; A might be good at drawing or talking. Everyone is good in different ways. Putting down others doesn't make you taller — it makes you narrower. Grown-ups know how to see others' strengths.'')}
\end{truyen}

\begin{baihoc}
Chê người khác không làm mình giỏi hơn; nhìn ra điểm mạnh của họ mới làm mình rộng ra.
\textit{(Putting others down doesn't make you better; seeing their strengths makes you broader.)}
\end{baihoc}

\begin{ghinhoanh}
\textbf{Short English to remember:}

Don't look down on others; everyone is good in different ways.

\end{ghinhoanh}

\begin{tuhocnhanh}
1. Em có hay so và chê bạn yếu hơn mình không? \textit{Do you often compare and put down those weaker than you?}

2. Điểm mạnh của vài bạn trong lớp em là gì? \textit{What are some classmates' strengths?}
\end{tuhocnhanh}
\ngancach

\section{Một lời nói tốt có thể thay đổi cả ngày của người khác.}
\begin{truyen}{Một lời nói tốt}{One Good Word}
\chuhoa{T}{hầy kể: ``Có em đi học buồn vì bị mắng ở nhà. Bạn bên cạnh nói: \textit{cậu hôm nay viết chữ đẹp đấy}. Chỉ một câu thôi. Em ấy kể lại với thầy: cả ngày em ấy đỡ hẳn. Một lời nói tốt có thể thay đổi cả ngày của người khác. Các em thử nói một câu tốt với ai đó mỗi ngày.''}
\textit{(The teacher said: ``A student came to school sad after being scolded at home. The one next to him said: \textit{your handwriting looks nice today}. Just one sentence. That student later said it made his whole day better. One good word can change someone's whole day. Try saying one good thing to someone every day.'')}
\end{truyen}

\begin{baihoc}
Lời nói tốt không tốn tiền nhưng giá trị rất lớn — với người nghe và với chính mình.
\textit{(Good words cost nothing but are worth a lot — to the listener and to yourself.)}
\end{baihoc}

\begin{ghinhoanh}
\textbf{Short English to remember:}

One kind word can change someone's whole day.

\end{ghinhoanh}

\begin{tuhocnhanh}
1. Hôm nay em đã nói một lời tốt với ai chưa? \textit{Did you say one kind word to someone today?}

2. Em nhớ lời tốt nào ai đó từng nói với em không? \textit{Do you remember a kind word someone said to you?}
\end{tuhocnhanh}
\ngancach

\section{Thái độ của em hôm nay sẽ quyết định bạn bè của em ngày mai.}
\begin{truyen}{Thái độ quyết định bạn bè}{Attitude Shapes Your Friends}
\chuhoa{T}{hầy nói: ``Em hay cáu, hay chê, hay đổ lỗi — thì bạn bè sẽ xa dần hoặc chỉ đến khi cần gì. Em hay lắng nghe, hay giúp đỡ, hay vui vẻ — thì bạn tốt sẽ tìm đến. Thái độ hôm nay quyết định bạn bè ngày mai. Em muốn xung quanh mình là người thế nào, hãy là người ấy trước.''}
\textit{(The teacher said: ``If you're often angry, critical, or blame others — friends will drift away or only show up when they need something. If you listen, help, and stay positive — good friends will find you. Today's attitude decides tomorrow's friends. If you want certain people around you, be that kind of person first.'')}
\end{truyen}

\begin{baihoc}
Bạn bè phản chiếu thái độ của mình; muốn bạn tốt thì hãy sống tốt trước.
\textit{(Friends reflect your attitude; want good friends, live well first.)}
\end{baihoc}

\begin{ghinhoanh}
\textbf{Short English to remember:}

Your attitude today shapes your friends tomorrow.

\end{ghinhoanh}

\begin{tuhocnhanh}
1. Thái độ của em với bạn bè đang thế nào? \textit{What is your attitude toward friends like?}

2. Em muốn bạn bè của mình là người thế nào? \textit{What kind of friends do you want?}
\end{tuhocnhanh}
\ngancach

\section{Khi em giúp người khác học tốt hơn, em cũng học tốt hơn.}
\begin{truyen}{Giúp bạn là giúp mình}{Helping Others Learn Helps You}
\chuhoa{T}{hầy khuyên: ``Khi em giảng lại cho bạn — dù bạn hỏi bài đơn giản — em đang tự ôn lại và sắp xếp lại ý trong đầu. Nhiều em giỏi lên rõ rệt sau khi hay kèm bạn. Giúp người khác học không mất thời gian; đó là cách em học sâu hơn.''}
\textit{(The teacher advised: ``When you explain to a friend — even a simple question — you're reviewing and organizing ideas in your head. Many students improve clearly after tutoring others. Helping others learn doesn't waste time; it's how you learn more deeply.'')}
\end{truyen}

\begin{baihoc}
Dạy lại người khác là cách ôn tốt nhất; ai hay giúp bạn học thường cũng học vững hơn.
\textit{(Teaching others is the best review; those who help others learn often learn more firmly themselves.)}
\end{baihoc}

\begin{ghinhoanh}
\textbf{Short English to remember:}

When you help others learn, you learn better too.

\end{ghinhoanh}

\begin{tuhocnhanh}
1. Em có hay giúp bạn hiểu bài không? \textit{Do you often help friends understand the lesson?}

2. Sau khi giảng cho bạn, em có thấy mình hiểu rõ hơn không? \textit{After explaining to a friend, do you feel you understand better?}
\end{tuhocnhanh}
\ngancach

\section{Người mạnh thật sự là người biết kiểm soát bản thân.}
\begin{truyen}{Kiểm soát bản thân}{Controlling Yourself}
\chuhoa{C}{ó em hay nổi nóng khi bị trêu. Thầy nói: ``Người mạnh không phải người đánh lại hay la to. Người mạnh là người biết dừng lại — không nói điều hối hận, không làm điều tổn thương. Kiểm soát được bản thân mới là sức mạnh thật. Em thử lần sau bị trêu: hít thở, đếm đến ba, rồi mới nói.''}
\textit{(A student often lost his temper when teased. The teacher said: ``The strong one isn't who hits back or shouts. The strong one is who can stop — not say what they'll regret, not do what hurts. Controlling yourself is real strength. Next time someone teases you: breathe, count to three, then speak.'')}
\end{truyen}

\begin{baihoc}
Sức mạnh thật là khi mình làm chủ được phản ứng của mình — không để cơn giận làm chủ.
\textit{(Real strength is when you master your own reaction — not let anger master you.)}
\end{baihoc}

\begin{ghinhoanh}
\textbf{Short English to remember:}

The truly strong person is the one who can control themselves.

\end{ghinhoanh}

\begin{tuhocnhanh}
1. Em có kiểm soát được mình khi bực không? \textit{Can you control yourself when you're upset?}

2. Làm sao để ``dừng lại'' trước khi nói hoặc làm điều hối hận? \textit{How can you pause before saying or doing something you'll regret?}
\end{tuhocnhanh}
\ngancach

\section{Trí tuệ làm em nổi bật, nhưng nhân cách làm em được tôn trọng.}
\begin{truyen}{Trí tuệ và nhân cách}{Intelligence and Character}
\chuhoa{T}{hầy nói: ``Trí tuệ khiến người ta chú ý em — điểm cao, trả lời hay. Nhưng nhân cách mới khiến người ta tôn trọng em — trung thực, giữ lời, không nói xấu, biết cảm ơn. Em có thể nổi bật nhờ trí tuệ; em được tôn trọng nhờ nhân cách. Cả hai mới đủ.''}
\textit{(The teacher said: ``Intelligence makes people notice you — good grades, good answers. But character makes them respect you — honesty, keeping your word, no gossip, gratitude. You can stand out by smarts; you earn respect by character. You need both.'')}
\end{truyen}

\begin{baihoc}
Nổi bật nhờ trí tuệ; được tôn trọng lâu dài nhờ nhân cách.
\textit{(You stand out by intelligence; you're respected for the long run by character.)}
\end{baihoc}

\begin{ghinhoanh}
\textbf{Short English to remember:}

Intelligence makes you stand out; character makes you respected.

\end{ghinhoanh}

\begin{tuhocnhanh}
1. Em đang chú ý rèn trí tuệ hay nhân cách? \textit{Are you focusing on intelligence or character?}

2. Nhân cách với em nghĩa là gì (cụ thể)? \textit{What does character mean to you in practice?}
\end{tuhocnhanh}
\ngancach

\section{Hãy học cách lắng nghe trước khi muốn người khác lắng nghe mình.}
\begin{truyen}{Lắng nghe trước}{Listen First}
\chuhoa{T}{hầy nhắc cả lớp: ``Nhiều em muốn nói, muốn được nghe, nhưng không chịu nghe người khác. Lắng nghe là cách tôn trọng và cũng là cách học — từ thầy, từ bạn, từ người lớn. Khi em biết lắng nghe, người ta mới muốn lắng nghe em. Hãy học lắng nghe trước.''}
\textit{(The teacher reminded the class: ``Many of you want to speak and be heard but won't listen to others. Listening is respect and also learning — from the teacher, from friends, from adults. When you know how to listen, people will want to listen to you. Learn to listen first.'')}
\end{truyen}

\begin{baihoc}
Muốn được nghe thì trước hết phải biết nghe; lắng nghe là nền tảng của giao tiếp và học tập.
\textit{(To be heard you must first know how to hear; listening is the base of communication and learning.)}
\end{baihoc}

\begin{ghinhoanh}
\textbf{Short English to remember:}

Learn to listen before you ask others to listen to you.

\end{ghinhoanh}

\begin{tuhocnhanh}
1. Em có hay ngắt lời hoặc không chú ý khi người khác nói không? \textit{Do you often interrupt or not pay attention when others speak?}

2. Lắng nghe giúp em học và quan hệ thế nào? \textit{How does listening help your learning and relationships?}
\end{tuhocnhanh}
\ngancach

\section{Một lớp học tốt là lớp học biết tôn trọng nhau.}
\begin{truyen}{Lớp biết tôn trọng nhau}{A Class That Respects Each Other}
\chuhoa{T}{hầy nói: ``Lớp học tốt không phải lớp toàn học sinh giỏi. Lớp học tốt là lớp không ai chê ai khi sai, không cười khi ai phát biểu sai, không nói xấu sau lưng. Mọi người tôn trọng nhau — thầy trò, bạn với bạn. Trong lớp như thế ai cũng dám hỏi, dám sai, dám tiến.''}
\textit{(The teacher said: ``A good class isn't one where everyone is top. A good class is where no one puts others down for mistakes, no one laughs when someone gives a wrong answer, no one gossips behind backs. Everyone respects each other — teacher and students, friend and friend. In such a class everyone dares to ask, to be wrong, to improve.'')}
\end{truyen}

\begin{baihoc}
Tôn trọng nhau tạo không khí an toàn; trong không khí ấy mới có học thật.
\textit{(Mutual respect creates a safe climate; only in that climate can real learning happen.)}
\end{baihoc}

\begin{ghinhoanh}
\textbf{Short English to remember:}

A good class is one that respects each other.

\end{ghinhoanh}

\begin{tuhocnhanh}
1. Lớp em có đang tôn trọng nhau không? \textit{Does your class respect each other?}

2. Em có thể làm gì để lớp tôn trọng nhau hơn? \textit{What can you do to make the class respect each other more?}
\end{tuhocnhanh}
\ngancach
